\chapter{Introduction}

\section{Terminology}

To avoid ambiguity, it is important to establish definitions for
terms used throughout this paper. The terms below may have other
broader definitions, but I will focus on certain narrow meanings
relevant to my work.

A \textbf{locale} is a specific set of conventions of a cultural region.
It is synonymous but not equivalent to \textit{language}.
For example, being a single language, the American and British English
locales have different norms.

\textbf{Localization} (l10n) is the process of adapting software
to different \textit{locales} and, possibly, adopting their
conventions for better integration into a cultural landscape
(e.g., \texttt{MM/DD/YYYY} vs. \texttt{DD.MM.YYYY} as the date format).

\textbf{Internationalization} (i18n) refers to designing software
in a way that it can support \textit{localization} without major
refactoring; for example, by keeping all natural language strings together.

These terms are reintroduced continuously in various sources.
\cite{barbour1997, he2002, couto2026}

\section{Brief Overview}

My work presents a standalone CLI utility for software developers
who want to internationalize and localize their applications.

The main unique features of my project are ecosystem-independence
(supporting multiple programming languages) and type-safety (through
a compilation step with code generation).

My project is shared under the MIT license and is publicly available
on GitHub: \url{https://github.com/oscarbash/mulan}.

\section{Problem Statement}

Many i18n frameworks exist. However, most of them are designed
to only support one ecosystem (e.g., JavaScript, Python, Android).
Different frameworks have different developer workflows and vary
by their quality and capabilities.

Although programming-language-agnostic tools exist, they are not
type-safe. Most frameworks depend on runtime lookups, which may be
error-prone and lead to poorer developer experience.

Currently, no single i18n solution offers ecosystem independence
with consistent ergonomics and type-safety for every supported platform.
I am looking to fill this gap with my work.

\section{Project Idea}

I present a new i18n framework that combines these features.
{\hypersetup{urlcolor=black}\href{https://xkcd.com/927/}{[927]}} This
framework is called \textbf{Mulan} (portmanteau of \textit{multi} +
\textit{language}). The goal of my tool is to consistently provide
the same functionality, type-safety, and developer experience for every
supported ecosystem.

\textit{Описать как я планирую подойти к проблеме. Какие стратегии
  существуют. Описать мой подход. Сказать, что в Literature Review
  будет описание альтернатив. Сказать чего я хочу добиться,
  acceptability criteria, как проект будет расширяться, какие языки и
функции будет поддерживать. Возможно, привести иллюстративные примеры, даты.}

CLDR \cite{cldr}, ICU \cite{icu}, ICU4X \cite{icu4x}

\section{Acceptance Criteria}

Rzk \cite{rzk}
...
